En los últimos años, la inteligencia artificial ha evolucionado hasta alcanzar la capacidad de generar, testear y ejecutar código con una competencia equiparable a la de un ingeniero cualificado. No obstante, continúa siendo una herramienta que requiere de un usuario capaz de guiarla y supervisar sus resultados. Un estudiante que se inicia en la programación carece de dichas facultades analíticas, sin embargo, tiene a su disposición una tecnología capaz de resolver, en cuestión de segundos, los retos iniciales inherentes al aprendizaje del desarrollo de software.

La prohibición del uso de estas herramientas en el ámbito académico puede resultar contraproducente, puesto que el aprendizaje de la programación puede generar altos niveles de frustración en principiantes. Esta situación conlleva el riesgo de que el alumnado abandone el estudio o recurra a soluciones automatizadas que impidan el desarrollo del pensamiento computacional, facultad indispensable para establecer los fundamentos técnicos que posteriormente permitirán aprovechar todo el potencial de estas tecnologías avanzadas de manera eficaz.

Ante este nuevo paradigma, el enfoque del programador profesional no debe residir exclusivamente en la sintaxis de los diversos lenguajes de programación. Por el contrario, el valor diferencial se desplaza hacia el diseño de sistemas, la descomposición modular en subtareas para su implementación y la selección estratégica de las herramientas y tecnologías más adecuadas para cada requerimiento técnico.

Bajo este contexto, SocratiCode nace como una respuesta técnica y pedagógica a la actual dicotomía entre la prohibición del uso de la IA y la automatización excesiva en el aprendizaje. El proyecto propone una arquitectura basada en modelos de lenguaje locales que, en lugar de generar soluciones directas, actúan como tutores socráticos. Al integrar un entorno de ejecución seguro con una asistencia iterativa, la plataforma busca transformar la IA en un catalizador del pensamiento computacional, garantizando que el alumno mantenga el control del flujo lógico y el esfuerzo cognitivo en su aprendizaje.